\chapter{Topological Prerequisites}

To fully leverage the power of measure theory in probability (especially constructing Borel \texorpdfstring{$\sigma$-}{sigma-}algebras and understanding convergence), we need a solid understanding of \textbf{Topology}.

A topology generalizes the notion of ``closeness'' and ``limits'' without necessarily needing a distance function (metric), although in finance we mostly work with metric spaces.

\section{Metric Spaces}

\begin{definition}[Metric Space]
A \textbf{metric space} is a pair $(M, d)$ where $M$ is a set and $d: M \times M \to [0, \infty)$ is a distance function satisfying:
\begin{enumerate}
    \item \textbf{Identity}: $d(x, y) = 0 \iff x = y$.
    \item \textbf{Symmetry}: $d(x, y) = d(y, x)$.
    \item \textbf{Triangle Inequality}: $d(x, z) \leq d(x, y) + d(y, z)$.
\end{enumerate}
\end{definition}

The most important building block in a metric space is the \textbf{Open Ball}.

\begin{definition}[Open Ball]
The open ball of radius $r > 0$ centered at $x \in M$ is:
\[
B(x, r) := \{ y \in M : d(x, y) < r \}.
\]
\end{definition}

In $\mathbb{R}$ with the Euclidean metric $d(x, y) = |x - y|$, the open ball is simply the open interval $(x-r, x+r)$. In $\mathbb{R}^2$, it is the interior of a disk.

\section{General Topology}

While metric spaces are intuitive, the rigorous definition of ``measurable events'' relies on the more general concept of \textbf{Open Sets}.

\begin{definition}[Topology]
A \textbf{topology} on a set $\Omega$ is a collection $\mathcal{T}$ of subsets of $\Omega$ (called \textbf{open sets}) such that:
\begin{enumerate}
    \item $\emptyset \in \mathcal{T}$ and $\Omega \in \mathcal{T}$.
    \item Any \textbf{union} (finite or infinite) of open sets is open.
    \item Any \textbf{finite intersection} of open sets is open.
\end{enumerate}
\end{definition}

In a metric space, a set $U$ is open if for every point $x \in U$, there exists some small ball entirely contained in $U$ ($B(x, \epsilon) \subset U$).

\subsection{Closed Sets}
A set $F$ is \textbf{closed} if its complement $F^c = \Omega \setminus F$ is open.
\textbf{Intuition}: Closed sets ``contain their boundaries''.
\begin{example}
In $\mathbb{R}$:
\begin{itemize}
    \item $(0, 1)$ is open.
    \item $[0, 1]$ is closed.
    \item $[0, 1)$ is neither open nor closed.
\end{itemize}
\end{example}

\subsection{Closure, Interior, Boundary}
For any set $A$:
\begin{itemize}
    \item \textbf{Interior} $A^\circ$: The largest open set contained in $A$.
    \item \textbf{Closure} $\bar{A}$: The smallest closed set containing $A$.
    \item \textbf{Boundary} $\partial A = \bar{A} \setminus A^\circ$.
\end{itemize}
Crucially, a set $A$ is closed if and only if $A = \bar{A}$.

\section{Continuity}

In calculus, we define continuity using $\epsilon-\delta$. In topology, we have a much cleaner definition that connects directly to $\sigma$-algebras.

\begin{definition}[Continuous Function]
A function $f: (X, \mathcal{T}_X) \to (Y, \mathcal{T}_Y)$ is \textbf{continuous} if the inverse image of every open set is open.
\[
\forall V \in \mathcal{T}_Y, \quad f^{-1}(V) \in \mathcal{T}_X.
\]
\end{definition}
This definition explains why random variables are defined as measurable functions: variables are to $\sigma$-algebras what continuous functions are to topologies.

\section{Important Properties for Probability}

\subsection{Separability}
A topological space is \textbf{separable} if it contains a countable dense subset.
\begin{itemize}
    \item $\mathbb{R}$ is separable because $\mathbb{Q}$ (rationals) is countable and dense in $\mathbb{R}$.
    \item \textbf{Why it matters}: Separability prevents the space from being ``too big''. It ensures that the Borel $\sigma$-algebra is well-behaved and that the support of a measure is unique.
\end{itemize}

\subsection{Completeness}
A metric space is \textbf{complete} if every Cauchy sequence converges to a limit within the space.
\begin{itemize}
    \item $\mathbb{R}$ is complete (Cauchy sequences converge to real numbers).
    \item $\mathbb{Q}$ is not complete (a sequence of rationals can converge to $\sqrt{2} \notin \mathbb{Q}$).
\end{itemize}

\subsection{Polish Spaces}
A \textbf{Polish Space} is a separable, completely metrizable topological space.
\begin{itemize}
    \item Examples: $\mathbb{R}$, $\mathbb{R}^d$, $C([0, T], \mathbb{R})$ (continuous functions, via Wiener measure).
    \item These are the ``nice'' spaces where modern probability theory lives.
\end{itemize}

\subsection{Compactness}
A set $K$ is \textbf{compact} if every open cover has a finite subcover.
In $\mathbb{R}^d$ (Heine-Borel theorem), a set is compact if and only if it is \textbf{closed} and \textbf{bounded}.
Compactness is crucial for proving the existence of limits (Prohorov's Theorem talks about ``tightness'', which is compactness for measures).
